\documentclass{smfbook}
\usepackage{sga1-smf}
%
%%%%%%%%%%%%%%%%%%%%%%%%%%%%%%%%%%%%%%%%%%%%%%%%%%%%%%%%%%%%%%%%%%%%%%%%
%
% This is a LaTeX 2e version of the book ``Rev{\^e}tements {\'E}tales et
% Groupe Fondamental'', Lecture Notes in Mathematics, 224,
% Springer-Verlag, Berlin-Heidelberg-New York, 1971, by Alexander
% Grothendieck et. al.
%
% Typesetting was done by volunteers, participating in a project on
% which one can find more details at
%
%                 http://www.math.leidenuniv.nl/~edix/
%
% and has been completed by the Societe mathematique de France.
%
% This file produces two versions (0 and 1, say), depending on the value
% of the boolean variable orig below. Version 0 (orig = true) should
% reproduce the original text, whereas version 1 (orig = false) takes
% reported errors into account and contains updating remarks by Michel Raynaud.
% The differences between the two versions
% are documented in the main TeX file. Both versions can be found in the
% archives `http://arxiv.org/'.
% Version 1 is also published by the Societe mathematique de France (SMF)
% in the series 'Documents Mathematiques', available at
%                 http://smf.emath.fr/
% 
%
% In order to compile this latex file, you need the files
% sga1-smf.tex (master file, this file),
% sga1-smf.sty (the specfic macro file),
% smf_doc-math_3_01.tex (the main TeX file, identical to the published file,
% input-ted below)
% the SMF standard classes
% smfbook.cls, smfenum.sty, smfthm.sty
% available at
%                 http://smf.emath.fr/Publications/Formats/
% and the standard AMS LaTeX packages.
% You can edit this file according to your own taste, and you
% adapt the boolean variable `orig' below as desired. Then you execute
% the following commands (in a unix type environment):
%
%   latex sga1-smf.tex (three times, say; the third time you should only get
%                   underfull hboxes and messages from xymatrix and hyperref);
%   makeindex sga1-smf
%   latex sga1-smf.tex (twice, maybe, to get the page numbers of the two
%                   indexes in the table of contents right).
%
% Then you have a correct sga1-smf.dvi dvi file (with hypertext facilities).
% Note that the Introduction starts on page vii, to keep a perfect
% correspondence with the new published edition.
%
% If you want to make a pdf file (with hypertext facilities) then you
% replace the five commands `latex' above by `pdflatex'. 
%
%
% Of course, you do not need to compile the latex file if you are
% happy with dvi, postscript, or pdf, as you can download those from 
% the arXiv archives. 
%
% If you find mistakes in this file, or have suggestions on how to
% improve it, please report them to Bas Edixhoven. If you pass on your
% own modified version of this file (and please think about the
% consequences before doing that), then add a text at its beginning
% and on the first page of the output, mentioning that it is a modified
% version. Do not modify this part of the file.
%
%%%%%%%%%%%%%%%%%%%%%%%%%%%%%%%%%%%%%%%%%%%%%%%%%%%%%%%%%%%%%%%%%%%%%%%

%% In order to make the two versions with a single file. 
%% With `\setboolean{orig}{true}' you get the original version, with
%% `\setboolean{orig}{false}' you get the corrected version. 

\setboolean{orig}{false}




\input{smf_doc-math_3_01.tex}